\documentclass[12pt,a4paper]{article}

\usepackage[margin=2.5cm]{geometry}
\usepackage{amsmath, amssymb}
\usepackage{graphicx}
\usepackage{hyperref}
\usepackage{listings}
\usepackage{xcolor}
\usepackage{booktabs}
\usepackage{array}
\usepackage{parskip}
\usepackage{titlesec}
\usepackage{fancyhdr}
\usepackage{microtype}

\hypersetup{
    colorlinks=true,
    linkcolor=blue!60!black,
    urlcolor=blue!60!black,
}

\lstset{
    language=C++,
    basicstyle=\ttfamily\small,
    keywordstyle=\color{blue!70!black}\bfseries,
    commentstyle=\color{green!50!black}\itshape,
    stringstyle=\color{red!60!black},
    numbers=left,
    numberstyle=\tiny\color{gray},
    numbersep=8pt,
    frame=single,
    framerule=0.4pt,
    rulecolor=\color{gray!40},
    breaklines=true,
    tabsize=4,
    showstringspaces=false,
    backgroundcolor=\color{gray!5},
}

\lstdefinestyle{bash}{
    language=bash,
    basicstyle=\ttfamily\small,
    keywordstyle=\color{blue!70!black}\bfseries,
    commentstyle=\color{green!50!black}\itshape,
    numbers=left,
    numberstyle=\tiny\color{gray},
    numbersep=8pt,
    frame=single,
    framerule=0.4pt,
    rulecolor=\color{gray!40},
    breaklines=true,
    backgroundcolor=\color{gray!5},
}

\pagestyle{fancy}
\fancyhf{}
\rhead{CCO-projects MD Code Report}
\lhead{\leftmark}
\cfoot{\thepage}
\renewcommand{\headrulewidth}{0.4pt}

\title{\textbf{Molecular Dynamics Simulation\\
in \texttt{CCOptimization}}\\[0.5em]
\large A Detailed Code Walkthrough}
\author{CCO-projects --- \texttt{generalized\_stealthy} branch}
\date{\today}

\begin{document}

\maketitle
\tableofcontents
\newpage

%=============================================================================
\section{Overview}
%=============================================================================

The \texttt{CCOptimization} module implements NVT molecular dynamics (MD)
simulations of point particles interacting through a \emph{stealthy
collective-coordinate potential}, optionally combined with a soft-core
repulsion. The goal is to sample finite-temperature configurations in which
the structure factor $S(\mathbf{k})$ is suppressed (or exactly zero at $T=0$)
over a prescribed union of annular shells in reciprocal space.

The MD pipeline spans three files:
\begin{itemize}
    \item \texttt{Potential/MD\_System.h} / \texttt{MD\_System.cpp} ---
          the integrator and thermostat (\texttt{ParticleMolecularDynamics} class).
    \item \texttt{CCOptimization/main.cpp} --- the CLI driver; reads parameters,
          builds the potential, and orchestrates the simulation stages.
    \item \texttt{CCOptimization/ex\_MD\_slurm.sh} --- the SLURM submission
          script; sets physical and computational parameters and calls the
          executable.
\end{itemize}

%=============================================================================
\section{The Stealthy Potential}
%=============================================================================

\subsection{Collective-coordinate energy}

The interaction energy is the collective-coordinate (CC) potential
\begin{equation}
    \Phi = \frac{1}{V} \sum_{\mathbf{k} \in \mathcal{C}}
           \tilde{V}(\mathbf{k})\,|\hat{\rho}(\mathbf{k})|^2
    \label{eq:CC}
\end{equation}
where $V$ is the volume of the periodic simulation cell,
$\mathcal{C}$ is the set of constrained wavevectors, and
\begin{equation}
    \hat{\rho}(\mathbf{k}) = \sum_{j=1}^{N} e^{i\mathbf{k}\cdot\mathbf{r}_j}
    = \rho_\text{real}(\mathbf{k}) + i\,\rho_\text{imag}(\mathbf{k})
\end{equation}
is the collective coordinate (structure factor amplitude).
The weight function $\tilde{V}(\mathbf{k})$ controls the shape of the
potential within the excluded region (flat, overlap, or power-law; set via
the \texttt{vtilde} argument).

\subsection{Constrained wavevector set --- generalized stealthy shells}

On the \texttt{generalized\_stealthy} branch, the excluded region is a
\emph{union of $M$ annular shells}:
\begin{equation}
    \mathcal{C} = \left\{\mathbf{k} :
    |\mathbf{k}| \in \bigcup_{n=1}^{M}
    \bigl[k_1^{(n)},\; k_1^{(n)} + \delta^{(n)}\bigr]\right\}
\end{equation}
Each shell $n$ is specified by a lower bound $k_1^{(n)}$ and a width
$\delta^{(n)}$, both in unit-number-density units.
The special case $M=1$, $k_1^{(1)}=0$ recovers the original single-shell
stealthy potential.

\subsection{Force on particle $i$}

Taking the gradient of Eq.~\eqref{eq:CC} with respect to $\mathbf{r}_i$:
\begin{equation}
    \mathbf{F}_i^{\text{CC}} = -\frac{2}{V}
    \sum_{\mathbf{k}\in\mathcal{C}} \tilde{V}(\mathbf{k})
    \Bigl[
        \rho_\text{imag}(\mathbf{k})\cos(\mathbf{k}\cdot\mathbf{r}_i)
        - \rho_\text{real}(\mathbf{k})\sin(\mathbf{k}\cdot\mathbf{r}_i)
    \Bigr]\mathbf{k}
    \label{eq:force}
\end{equation}

\subsection{Soft-core repulsion}

When \texttt{vareps0} $> 0$, a contact (soft-core) potential is added:
\begin{equation}
    \Phi_\text{rep} = \varepsilon_0 \sum_{i<j}
    \left(1 - \frac{r_{ij}}{\sigma_i + \sigma_j}\right)^2
    \Theta(\sigma_i + \sigma_j - r_{ij})
\end{equation}
where $\sigma_i$ is the radius of particle $i$ and $\Theta$ is the
Heaviside step function.
The total energy is $\Phi_\text{tot} = u_0\,\Phi_\text{CC} + e_0\,\Phi_\text{rep}$.

%=============================================================================
\section{The \texttt{ParticleMolecularDynamics} Class}
%=============================================================================

Defined in \texttt{Potential/MD\_System.h}, this class owns the full
dynamical state of the system.

\subsection{Member variables}

\begin{center}
\begin{tabular}{lll}
\toprule
Variable & Type & Description \\
\midrule
\texttt{Position} & \texttt{Configuration} & Particle positions (relative coords) \\
\texttt{Velocity} & \texttt{vector<GeometryVector>} & Velocities, one per particle \\
\texttt{Acceleration} & \texttt{vector<GeometryVector>} & Accelerations (= forces / mass) \\
\texttt{Mass} & \texttt{vector<double>} & Particle masses (all set to 1.0) \\
\texttt{TimeStep} & \texttt{double} & MD time step $\Delta t$ \\
\texttt{Time} & \texttt{double} & Cumulative simulation time \\
\bottomrule
\end{tabular}
\end{center}

\subsection{Acceleration computation --- \texttt{\_GetAcc}}

\begin{lstlisting}
void ParticleMolecularDynamics::_GetAcc(Potential & potential)
{
    potential.AllForce(this->Acceleration);   // F_i for all i at once
    for(int i=0; i<this->Acceleration.size(); i++)
        this->Acceleration[i].MultiplyFrom(1.0 / this->Mass[i]);
}
\end{lstlisting}

\texttt{AllForce} calls \texttt{Force(i)} for every particle $i$, which
evaluates Eq.~\eqref{eq:force}. Crucially, $\rho(\mathbf{k})$ is computed
once for the full configuration before any force is evaluated (via the
\texttt{RhoValid} dirty-bit cache).

%=============================================================================
\section{The Velocity-Verlet Integrator --- \texttt{Evolve}}
%=============================================================================

\texttt{Evolve(n, potential)} advances the system by $n$ time steps of
size $\Delta t$ using the standard velocity-Verlet algorithm.
\textbf{All $N$ particles are updated at every step.}

\subsection{Algorithm (one step)}

\begin{align}
    \mathbf{v}_i\!\left(t + \tfrac{\Delta t}{2}\right)
        &= \mathbf{v}_i(t) + \tfrac{\Delta t}{2}\,\mathbf{a}_i(t)
        && \text{(half-kick, all } i\text{)}
    \\
    \mathbf{r}_i(t + \Delta t)
        &= \mathbf{r}_i(t) + \Delta t\,\mathbf{v}_i\!\left(t + \tfrac{\Delta t}{2}\right)
        && \text{(drift, all } i\text{)}
    \\
    \mathbf{a}_i(t + \Delta t)
        &= \mathbf{F}_i / m
        && \text{(force eval, all } i \text{ globally)}
    \\
    \mathbf{v}_i(t + \Delta t)
        &= \mathbf{v}_i\!\left(t + \tfrac{\Delta t}{2}\right)
           + \tfrac{\Delta t}{2}\,\mathbf{a}_i(t+\Delta t)
        && \text{(half-kick, all } i\text{)}
\end{align}

\subsection{Code}
\begin{lstlisting}
void ParticleMolecularDynamics::Evolve(size_t repeat, Potential & potential)
{
    for(size_t u=0; u<repeat; u++)
    {
        // half-kick
        for(long i=0; i<this->NumParticle; i++)
            this->Velocity[i].AddFrom(0.5*this->TimeStep*this->Acceleration[i]);

        // drift
        for(long i=0; i<this->NumParticle; i++){
            GeometryVector newCart = this->Position.GetCartesianCoordinates(i)
                                   + this->TimeStep*this->Velocity[i];
            this->Position.MoveParticle(i,
                this->Position.CartesianCoord2RelativeCoord(newCart));
        }

        // force re-evaluation
        potential.SetConfiguration(this->Position);
        this->_GetAcc(potential);

        // half-kick
        for(long i=0; i<this->NumParticle; i++)
            this->Velocity[i].AddFrom(0.5*this->TimeStep*this->Acceleration[i]);

        this->Time += this->TimeStep;
    }
}
\end{lstlisting}

Velocity Verlet is symplectic and time-reversible, making it suitable for
MD. Energy conservation (in the microcanonical NVE sense) depends on the
choice of $\Delta t$.

%=============================================================================
\section{The Anderson Thermostat --- \texttt{AndersonEvolve}}
%=============================================================================

\subsection{Purpose}

\texttt{AndersonEvolve} implements the \textbf{Anderson thermostat} to
maintain NVT (constant temperature) dynamics. It augments each velocity-Verlet
step with stochastic velocity resets that couple the system to a fictitious
heat bath at temperature $T$.

\subsection{Mechanism}

After each complete velocity-Verlet step, a running accumulator
\texttt{ResetCount} is incremented by \texttt{ResetSpeedProbability} (= 0.01).
Whenever \texttt{ResetCount} $\geq 1$, one randomly chosen particle has
its velocity completely redrawn from the Maxwell-Boltzmann distribution
at temperature $T$:

\begin{lstlisting}
ResetCount += ResetSpeedProbability;      // 0.01 per step
for (; ResetCount > 0; ResetCount -= 1.0)
{
    size_t i = floor(gen.RandomDouble() * NumParticle);  // random particle
    double sigma = sqrt(kT / Mass[i]);
    for (DimensionType j = 0; j < Dimension; j++)
    {
        // Box-Muller: Gaussian with std = sigma
        this->Velocity[i].x[j] =
            sqrt(-2*log(gen.RandomDouble()))
            * sin(2*pi*gen.RandomDouble())
            * sigma;
    }
}
\end{lstlisting}

\subsection{Box-Muller transform}

The new velocity component is drawn from a Gaussian via the Box-Muller
transform. If $u_1, u_2 \sim \mathcal{U}(0,1)$, then
\begin{equation}
    v_j = \sqrt{-2\ln u_1}\;\sin(2\pi u_2)\;\sigma,
    \qquad \sigma = \sqrt{k_BT/m}
\end{equation}
is Gaussian with mean 0 and variance $k_BT/m$, which is the Maxwell-Boltzmann
distribution for a single velocity component:
\begin{equation}
    P(v_j) \propto \exp\!\left(-\frac{m v_j^2}{2k_BT}\right)
\end{equation}

\subsection{Kick rate}

With \texttt{ResetSpeedProbability} = 0.01:
\begin{itemize}
    \item One velocity reset occurs every \textbf{100 steps} (deterministic,
          via the accumulator).
    \item The chosen particle is \textbf{random} each time.
    \item For $N = 100$ particles, each particle is reset on average once
          every $100 \times 100 = 10{,}000$ steps.
\end{itemize}

The thermostat acts weakly: most dynamics between kicks are purely Newtonian,
preserving realistic trajectories while maintaining the target temperature
over long times.

%=============================================================================
\section{Velocity Initialization --- \texttt{SetRandomSpeed}}
%=============================================================================

\begin{lstlisting}
void ParticleMolecularDynamics::SetRandomSpeed(double kT, RandomGenerator & gen)
{
    for(size_t i=0; i<this->NumParticle; i++)
    {
        double sigma = sqrt(kT / this->Mass[i]);
        for (DimensionType j = 0; j < this->Dimension; j++)
            this->Velocity[i].x[j] =
                sqrt(-2*log(gen.RandomDouble()))
                * sin(2*pi*gen.RandomDouble()) * sigma;
    }
}
\end{lstlisting}

This is an \emph{instantaneous} reset of \textbf{all} $N \times d$ velocity
components simultaneously, drawn from the Maxwell-Boltzmann distribution
at temperature $T$. It is used during equilibration (Stage 1) to aggressively
drive the system from $T=0$ to the target temperature.

%=============================================================================
\section{The Simulation Pipeline --- \texttt{CollectiveCoordinateMD}}
%=============================================================================

The full MD run is managed by \texttt{CollectiveCoordinateMD} in
\texttt{CCOptimization/main.cpp}. It proceeds through three phases.

\subsection{Phase 0: Initial quench}

\begin{lstlisting}
RelaxStructure_NLOPT(*pConfig, *pPotential, 0.0, 0, 0.0, 1000);
psystem = new ParticleMolecularDynamics(*pConfig, TimeStep, 1.0);
\end{lstlisting}

The initial configuration (random or loaded) is first minimized to a
local energy minimum via LBFGS (\texttt{RelaxStructure\_NLOPT}).
This ensures the MD starts from a ground state at $T=0$ rather than
a high-energy random configuration. The \texttt{ParticleMolecularDynamics}
object is then created with all velocities set to zero.

\subsection{Phase 1: Equilibration (Stage 1)}

Runs \texttt{EquilibrateSamples} blocks (set by \texttt{numequilibration}).
Each block:
\begin{enumerate}
    \item \texttt{SetRandomSpeed(T)} --- redraw \emph{all} velocities from
          Maxwell-Boltzmann at $T$. This is a complete velocity reset,
          not the Anderson thermostat.
    \item \texttt{Evolve(StepPerSample/2)} --- pure NVE for $n/2$ steps.
    \item Repeat steps 1--2 once more.
\end{enumerate}

Total steps in equilibration: $\texttt{numequilibration} \times \texttt{samplesteps}$.
\textbf{No configurations are saved.}

The purpose is to heat the system rapidly from $T=0$ to the target
temperature. The brute-force velocity resets are intentionally aggressive ---
they are not physically realistic but quickly establish the correct
kinetic energy scale.

\subsection{Phase 2: Sampling (Stage 2)}

Runs \texttt{SampleNumber} = \texttt{Nc} blocks.
Each block:
\begin{enumerate}
    \item \texttt{AndersonEvolve(StepPerSample/2)} --- velocity Verlet
          with Anderson thermostat for $n/2$ steps.
    \item \texttt{Evolve(StepPerSample/2)} --- pure NVE for $n/2$ steps.
    \item \textbf{Save} the current configuration to the output
          \texttt{ConfigPack}.
\end{enumerate}

One configuration is saved per block, separated by
$\texttt{samplesteps} \times \Delta t$ time units.

\subsection{Timeline summary}

\begin{center}
\renewcommand{\arraystretch}{1.4}
\begin{tabular}{lccc}
\toprule
Phase & Blocks & Steps/block & Saved? \\
\midrule
Quench (LBFGS) & 1 & up to 1000 evals & No \\
Equilibration & \texttt{numequilibration} & \texttt{samplesteps} & No \\
Sampling & \texttt{Nc} & \texttt{samplesteps} & Yes (1 per block) \\
\bottomrule
\end{tabular}
\end{center}

With the example script parameters (\texttt{numequilibration}=10,
\texttt{samplesteps}=1000, \texttt{Nc}=20, $\Delta t$=0.05):

\begin{center}
\renewcommand{\arraystretch}{1.4}
\begin{tabular}{lrr}
\toprule
Phase & Total steps & Total time units \\
\midrule
Equilibration & 10,000 & 500 \\
Sampling & 20,000 & 1000 \\
\textbf{Total} & \textbf{30,000} & \textbf{1500} \\
\bottomrule
\end{tabular}
\end{center}

%=============================================================================
\section{Checkpoint and Restart --- \texttt{restore}}
%=============================================================================

When \texttt{restore} is passed, the code saves a binary checkpoint
\texttt{fname\_MD.MDDump} at the end (or when the wall-clock limit is reached).
The checkpoint stores the full dynamical state:
position, velocity, acceleration, time, stage counter, and step counter.

On the next invocation (same command with \texttt{restore}), the code
detects the checkpoint file, loads it, and \textbf{continues the trajectory
from where it stopped} --- no re-equilibration, no loss of progress.

This is essential for long cluster jobs. The wall-clock time limit is checked
at the start of every block:
\begin{lstlisting}
if (std::time(nullptr) > TimeLimit || std::time(nullptr) > ::TimeLimit)
{
    // save .MDDump checkpoint
    // write fname_continue.txt warning
    return 0;
}
\end{lstlisting}

%=============================================================================
\section{Stopping Criterion}
%=============================================================================

The code has \textbf{no automatic steady-state detection}. It stops when:
\begin{enumerate}
    \item All \texttt{Nc} sampling blocks are completed, or
    \item The wall-clock time limit (\texttt{timelimit} hours) is reached.
\end{enumerate}

Whether the system has truly equilibrated is the user's responsibility.
The log prints $E_p$ (potential energy) and $E_k$ (kinetic energy) at
each saved snapshot:
\begin{lstlisting}
std::cout << "2:" << i << "/" << SampleNumber
          << " \tE_relax=" << Ep
          << " \tE_k=" << Ek << '\n';
\end{lstlisting}

The user should inspect these values and verify that $E_p$ has stabilized
before using the saved configurations for analysis.

%=============================================================================
\section{The SLURM Script --- \texttt{ex\_MD\_slurm.sh}}
%=============================================================================

\subsection{Parameter encoding}

Physical k-space parameters are specified at unit number density and
converted to absolute units via:
\begin{equation}
    a = \left(\frac{\phi_\text{fake}}{V_d}\right)^{1/d}, \qquad
    K_{1,a}^{(n)} = k_1^{(n)} \cdot a, \qquad
    \delta_a^{(n)} = \delta^{(n)} \cdot a
\end{equation}
where $V_d$ is the volume of a unit hypersphere in $d$ dimensions.
The code divides back by $a$ to recover absolute values.
This encoding makes the script parameters independent of $N$ and box size.

\subsection{Shell string construction}

\begin{lstlisting}[style=bash]
shell_str="${M}"
for (( n=0; n<M; n++ )); do
    K1a=`echo "${K1s[$n]}*${a}"       | bc -l`
    deltaa=`echo "${deltas[$n]}*${a}" | bc -l`
    shell_str="${shell_str} ${K1a} ${deltaa}"
done
\end{lstlisting}

\subsection{CLI call}

\begin{lstlisting}[style=bash]
${exe} ${timelimit} ${seed} ${beg_idx} ${Verbosity} <<< \
    "${d} ${shell_str} $S0 $vareps0 ${sigma} ${phi_fake} \
     ${threads} ${N} ${Nc} random ${fname}_MD \
     MD $T $dt $int_samp $num_equi $restore run"
\end{lstlisting}

\subsection{Key parameters}

\begin{center}
\renewcommand{\arraystretch}{1.4}
\begin{tabular}{lll}
\toprule
Script variable & Default & Description \\
\midrule
\texttt{d} & 2 & Spatial dimension \\
\texttt{M} & 2 & Number of stealthy shells \\
\texttt{K1s} & (0.0, 0.5) & Lower bounds per shell (unit density) \\
\texttt{deltas} & (0.3, 0.2) & Widths per shell (unit density) \\
\texttt{S0} & 0.0 & Target $S(\mathbf{k})$; 0 = stealthy \\
\texttt{vareps0} & 1.0 & Soft-core repulsion strength \\
\texttt{N} & 100 & Number of particles \\
\texttt{Nc} & 20 & Configurations to collect \\
\texttt{temperature} & $10^{-3}$ & Dimensionless temperature (units of $v_0$) \\
\texttt{timestep} & 0.05 & MD time step $\Delta t$ \\
\texttt{samplesteps} & 1000 & Steps between saved snapshots \\
\texttt{numequilibration} & 10 & Equilibration blocks \\
\texttt{timelimit} & 4 & Wall-clock limit (hours) \\
\bottomrule
\end{tabular}
\end{center}

%=============================================================================
\section{Data Flow Diagram}
%=============================================================================

\begin{center}
\renewcommand{\arraystretch}{1.6}
\begin{tabular}{cl}
\toprule
Step & Action \\
\midrule
1 & \texttt{ex\_MD\_slurm.sh}: encode $k$-space shells, build CLI string \\
2 & \texttt{main.cpp / GetCCO}: parse stdin, build constraint list $\mathcal{C}$ \\
3 & \texttt{RepulsiveCCPotential}: store $({\bf k}, \tilde{V}({\bf k}))$ for each ${\bf k} \in \mathcal{C}$ \\
4 & \texttt{CollectiveCoordinateMD}: initial LBFGS quench \\
5 & Stage 1: \texttt{SetRandomSpeed} + \texttt{Evolve} $\times$ \texttt{numequilibration} \\
6 & Stage 2: \texttt{AndersonEvolve} + \texttt{Evolve} $\times$ \texttt{Nc}, save each \\
7 & Output: \texttt{fname\_MD.ConfigPack} with \texttt{Nc} configurations \\
\bottomrule
\end{tabular}
\end{center}

%=============================================================================
\section{Summary}
%=============================================================================

\begin{itemize}
    \item The integrator is \textbf{velocity Verlet} --- symplectic,
          time-reversible, applied to \emph{all} $N$ particles at every step.
    \item The thermostat is \textbf{Anderson}: one randomly chosen particle
          has its velocity redrawn from Maxwell-Boltzmann every 100 steps.
          Between kicks, dynamics are purely Newtonian (NVE).
    \item \textbf{Equilibration} uses brute-force full velocity resets
          (\texttt{SetRandomSpeed}) to heat the system from $T=0$ rapidly.
          No configs are saved.
    \item \textbf{Sampling} switches to the gentle Anderson thermostat.
          One configuration is saved every \texttt{samplesteps} $\times \Delta t$
          time units.
    \item There is \textbf{no automatic equilibration check}. The user must
          inspect $E_p$ convergence in the log.
    \item \textbf{Checkpointing} via \texttt{restore} allows seamless
          continuation across SLURM job submissions.
    \item The \texttt{generalized\_stealthy} branch supports $M$ annular
          shells in $k$-space; the script builds the shell string automatically.
\end{itemize}

\end{document}
